\documentclass[tikz, border=10pt]{standalone}
\usepackage{tikz}
\usetikzlibrary{shapes.geometric, arrows.meta, positioning, fit, backgrounds}

\tikzset{
  box/.style={
    rectangle, rounded corners=4pt, draw, text centered,
    minimum width=3.8cm, minimum height=0.8cm,
    font=\small, inner sep=6pt
  },
  decision/.style={
    diamond, draw, text centered, aspect=2,
    font=\small, inner sep=3pt
  },
  arrow/.style={-{Stealth[length=6pt]}, thick},
  darrow/.style={-{Stealth[length=6pt]}, thick, dashed},
  %--- colour palette ---
  config/.style   ={box, fill=gray!15,       draw=gray!60},
  setup/.style    ={box, fill=yellow!25,     draw=orange!60},
  batch/.style    ={box, fill=white,         draw=black!50},
  llm/.style      ={box, fill=blue!10,       draw=blue!50},
  parse/.style    ={box, fill=green!12,      draw=green!50!black},
  out/.style      ={box, fill=purple!10,     draw=purple!50},
  grouplabel/.style={font=\small\bfseries, anchor=north west},
}

\begin{document}
\begin{tikzpicture}[node distance=0.55cm and 1.4cm]

%% ---------------------------------------------------------------
%% CONFIG column (left)
%% ---------------------------------------------------------------
\node[config] (c1) {Load API key \& lab number};
\node[config, below=of c1] (c2) {Resolve paths:\\rubric / solution /\\starter / students dir};
\draw[arrow] (c1) -- (c2);

%% ---------------------------------------------------------------
%% SETUP column (far left, R-only)
%% ---------------------------------------------------------------
\node[setup, left=2.0cm of c1, yshift=-0.6cm] (s1) {Render .qmd $\to$ Markdown};
\node[setup, below=of s1] (s2) {Upload rubric, solution,\\starter to Files API};
\node[setup, below=of s2] (s3) {Create persistent Assistant\\(file\_search tool)};
\node[setup, below=of s3] (s4) {Save IDs to\\assistant\_config.json};
\draw[arrow] (s1) -- (s2);
\draw[arrow] (s2) -- (s3);
\draw[arrow] (s3) -- (s4);

%% ---------------------------------------------------------------
%% BATCH LOOP — assemble context
%% ---------------------------------------------------------------
\node[batch, right=2.0cm of c1, yshift=-0.6cm] (b1) {Discover student\\submission files};

\node[config, below=of b1]                     (ctx1) {Rubric};
\node[config, below=of ctx1]                   (ctx2) {Instructor solution};
\node[config, below=of ctx2]                   (ctx3) {Starter template};
\node[config, below=of ctx3]                   (ctx4) {Student submission};

\draw[arrow] (b1)   -- (ctx1);
\draw[arrow] (ctx1) -- (ctx2);
\draw[arrow] (ctx2) -- (ctx3);
\draw[arrow] (ctx3) -- (ctx4);

%% ---------------------------------------------------------------
%% LLM interaction (two parallel paths)
%% ---------------------------------------------------------------
\node[llm, right=2.0cm of ctx2, yshift=0.4cm]  (lpy) {\textbf{Python:} inline all context\\in Chat Completions\\(prompt caching)};
\node[llm, below=0.9cm of lpy]                 (lr)  {\textbf{R:} attach file IDs to thread;\\Assistants API retrieves\\files via file\_search};
\node[llm, below=0.9cm of lr]                  (ljson){Enforce JSON\\response format};
\node[llm, below=of ljson]                     (lpoll){Await response\\(R polls; Python synchronous)};

\draw[arrow] (ljson) -- (lpoll);

%% ctx4 fans out to both LLM paths
\draw[arrow] (ctx4.east) -- ++(0.5,0) |- (lpy.west);
\draw[arrow] (ctx4.east) -- ++(0.5,0) |- (lr.west);
\draw[arrow] (lpy.east)  -- ++(0.4,0) |- (ljson.west);
\draw[arrow] (lr.east)   -- ++(0.4,0) |- (ljson.west);

%% ---------------------------------------------------------------
%% PARSE & COLLECT
%% ---------------------------------------------------------------
\node[parse, right=2.0cm of ljson] (p1) {Extract JSON:\\questions: \{Q1…Qn:\\\ \ grade, feedback\}\\total, overall\_comment};
\node[parse, below=of p1]          (p2) {Flatten to one row\\per student};

\draw[arrow] (lpoll.east) -| (p1.south);
\draw[arrow] (p1) -- (p2);

%% ---------------------------------------------------------------
%% OUTPUT
%% ---------------------------------------------------------------
\node[out, below=of p2] (o1) {Append row to results};
\node[out, below=of o1] (o2) {Write grades CSV};

\draw[arrow] (p2) -- (o1);
\draw[arrow] (o1) -- (o2);

%% loop-back arrow
\draw[darrow] (o1.west) -- ++(-0.6,0) -- ++(0,5.2) -- (b1.west);

%% ---------------------------------------------------------------
%% CONFIG -> BATCH entry
%% ---------------------------------------------------------------
\draw[arrow] (c2.east) -- (b1.west);

%% R setup feeds into batch (dashed)
\draw[darrow] (s4.east) -- ++(0.5,0) |- node[font=\tiny, above, pos=0.3]
              {R: load saved IDs} (b1.west);

%% ---------------------------------------------------------------
%% GROUP BOXES
%% ---------------------------------------------------------------
\begin{scope}[on background layer]
  % Setup (R only)
  \node[draw=orange!60, fill=yellow!10, rounded corners=6pt,
        fit=(s1)(s2)(s3)(s4), inner sep=8pt] (setupbox) {};
  \node[grouplabel, text=orange!80!black] at (setupbox.north west)
        {Setup (R only)};

  % Config
  \node[draw=gray!60, fill=gray!8, rounded corners=6pt,
        fit=(c1)(c2), inner sep=8pt] (configbox) {};
  \node[grouplabel, text=gray!70!black] at (configbox.north west)
        {Configuration};

  % Batch loop
  \node[draw=black!40, fill=white, rounded corners=6pt,
        fit=(b1)(ctx1)(ctx2)(ctx3)(ctx4), inner sep=8pt] (batchbox) {};
  \node[grouplabel] at (batchbox.north west) {Batch loop (per student)};

  % LLM
  \node[draw=blue!50, fill=blue!5, rounded corners=6pt,
        fit=(lpy)(lr)(ljson)(lpoll), inner sep=8pt] (llmbox) {};
  \node[grouplabel, text=blue!60!black] at (llmbox.north west)
        {LLM Interaction};

  % Parse
  \node[draw=green!50!black, fill=green!5, rounded corners=6pt,
        fit=(p1)(p2), inner sep=8pt] (parsebox) {};
  \node[grouplabel, text=green!40!black] at (parsebox.north west)
        {Parse \& Collect};

  % Output
  \node[draw=purple!50, fill=purple!5, rounded corners=6pt,
        fit=(o1)(o2), inner sep=8pt] (outbox) {};
  \node[grouplabel, text=purple!60!black] at (outbox.north west)
        {Output};
\end{scope}

\end{tikzpicture}
\end{document}
